\documentclass[portuguese]{sbc2025}

\usepackage[misc,geometry]{ifsym}
\usepackage{aas_macros}
\usepackage[bottom]{footmisc}
\usepackage{tabularray}
\usepackage{afterpage}
\usepackage{url}
\usepackage{pifont}
\usepackage{amsmath}
\usepackage{listings}
\usepackage{booktabs}
\usepackage{multirow}

\setcitestyle{square}

\definecolor{engtitle}{rgb}{0.5,0.5,0.5}
\definecolor{orcidlogo}{rgb}{0.37,0.48,0.13}
\definecolor{darkblue}{rgb}{0.0,0.0,0.0}
\definecolor{codegray}{rgb}{0.95,0.95,0.95}
\hypersetup{colorlinks,breaklinks,
            linkcolor=darkblue,urlcolor=darkblue,
            anchorcolor=darkblue,citecolor=darkblue}

\lstset{
  backgroundcolor=\color{codegray},
  basicstyle=\small\ttfamily,
  breaklines=true,
  frame=single,
  rulecolor=\color{gray}
}

\jid{AM-UFSJ}
\jtitle{Aprendizado de Máquina, UFSJ, 2026/1}
\issn{}
\doi{}
\copyrightstatement{}
\jyear{2026}

\category{Trabalho Prático}

\title[Framework para Busca de Jogadores de Futebol]{%
  Framework para Busca de Jogadores de Futebol}

\engtitle{\textcolor{engtitle}{A Framework for Football Player Search}}

\author[Lima et al. 2026]{
\affil{\textbf{André Chagas Lima}~~[{Universidade Federal de São João del-Rei}~|%
  \href{mailto:andrechalima@aluno.ufsj.edu.br}{\textit{ andrechalima@aluno.ufsj.edu.br}}~]}

\affil{\textbf{Lucas Eduardo Bernardes de Paula}~~[{Universidade Federal de São João del-Rei}~|%
  \href{mailto:xxlucas0404xx@aluno.ufsj.edu.br}{\textit{ xxlucas0404xx@aluno.ufsj.edu.br}}~]}

\affil{\textbf{Messias Feres Curi Melo}~~[{Universidade Federal de São João del-Rei}~|%
  \href{mailto:messiasferes127@aluno.ufsj.edu.br}{\textit{ messiasferes127@aluno.ufsj.edu.br}}~]}
}

% ============================================================
\begin{document}
\begin{frontmatter}

\maketitle

\begin{mail}
Departamento de Ciência da Computação, Universidade Federal de São João del-Rei,
Praça Frei Orlando, 170, Centro, São João del-Rei, MG, 36307-352, Brasil.
\end{mail}

% ---- Resumo em português ----
\begin{abstract-pt}
Este trabalho propõe o desenvolvimento de um \textit{framework} modular para busca
de jogadores de futebol a partir de estatísticas de desempenho da temporada 2024--2025.
O problema consiste em recuperar atletas compatíveis com perfis técnicos específicos
a partir de consultas em linguagem natural, apoiando tarefas de scouting, análise de
mercado e comparação de atletas. O \textit{framework} compara duas estratégias de
representação textual, a concatenação de \textit{labels} e a geração via modelo de
linguagem local (Llama 3.2), sobre dois métodos de recuperação: BM25 e Bi-encoder.
Como resultados, o sistema conta com pré-processamento completo, as duas
representações textuais, dois mecanismos de \textit{ground truth} (baseado em LLM e
baseado em regras ponderadas) e ambos os motores de recuperação operacionais, com
avaliação por Precision@$k$, Recall@$k$ e NDCG@$k$. A arquitetura modular, com
controle via \textit{ground truth}, permite experimentos reprodutíveis e comparações
sistemáticas no contexto de scouting.
\end{abstract-pt}

% ---- Abstract in English ----
\begin{abstract-en}
This work proposes a modular framework for football player search based on performance
statistics from the 2024--2025 season. The problem consists of retrieving athletes
compatible with specific technical profiles from natural language queries, supporting
scouting, market analysis, and player comparison tasks. The framework compares two
textual representation strategies, label concatenation and text generated by a local
language model (Llama 3.2), over two retrieval methods: BM25 and Bi-encoder. The
system features complete preprocessing, both textual representations, two ground truth
mechanisms (LLM-based and rule-based), and both operational retrieval engines,
evaluated with Precision@$k$, Recall@$k$, and NDCG@$k$. The modular architecture,
with ground truth control, enables reproducible experiments and systematic comparisons
in a sports scouting context.
\end{abstract-en}

\begin{pchaves}
Aprendizado de máquina; Scouting; Busca de jogadores; Recuperação de informação;
BM25; Bi-encoder; Modelos de linguagem; Ground Truth
\end{pchaves}

\begin{keywords}
Machine learning; Scouting; Player search; Information retrieval;
BM25; Bi-encoder; Language models; Ground Truth
\end{keywords}

\begin{dates}
\noindent{\sffamily\textbf{Entregue/Submitted:}} Maio 2026~~~$\bullet$~~~
{\sffamily\textbf{Disciplina:}} Aprendizado de Máquina~~~$\bullet$~~~
{\sffamily\textbf{Semestre:}} 2026/1
\end{dates}

\end{frontmatter}

% ============================================================
\section{Introdução}
\label{sec:intro}

O futebol é um dos esportes mais populares do mundo, movimentando bilhões de dólares
anualmente no mercado de transferências e atraindo investimentos crescentes em
tecnologias de análise de desempenho \citep{delatorre2021}. Identificar, de forma
automatizada, atletas que se encaixem em contextos específicos (seja para substituir
um jogador lesionado, reforçar uma posição ou apoiar decisões de mercado) é uma
tarefa de alta relevância para clubes, comissões técnicas e agências esportivas.

O processo tradicional de \textit{scouting} enfrenta quatro limitações centrais que
motivam este trabalho. Primeiro, a \textbf{lentidão da busca manual}: um analista
humano não consegue acompanhar simultaneamente milhares de atletas distribuídos em
dezenas de ligas. Segundo, a \textbf{dificuldade de leitura dos dados tabulares}: as
estatísticas brutas de desempenho são difíceis de interpretar diretamente, dificultando
a formação de uma visão holística do perfil do jogador. Terceiro, a
\textbf{similaridade entre atletas pode ser explorada automaticamente}: métricas de
distância e similaridade permitem identificar padrões e perfis equivalentes sem
depender exclusivamente do julgamento subjetivo do observador. Quarto, a
\textbf{necessidade de interpretabilidade}: a solução deve ser compreensível e útil
na prática, permitindo que o analista entenda e justifique as recomendações produzidas.

Nesse contexto, técnicas de recuperação de informação (\textit{Information Retrieval},
IR) oferecem ferramentas para estruturar e recuperar informações de forma automatizada
e escalável. O BM25 \citep{robertson2009} é o modelo probabilístico de recuperação
mais consolidado na literatura, servindo como \textit{baseline} forte em benchmarks
como o BEIR \citep{thakur2021}. Mais recentemente, Bi-encoders baseados em
Transformers, como o DPR \citep{karpukhin2020}, abriram novas possibilidades para
recuperar documentos a partir de consultas em linguagem natural, aproximando o sistema
de uma interface mais intuitiva para analistas esportivos.

Apesar do crescente interesse na área, os sistemas propostos na literatura
frequentemente adotam apenas uma estratégia de representação e um único método de
recuperação, dificultando comparações sistemáticas \citep{springer2026,researchgate2025}.
Há, portanto, uma lacuna para \textit{frameworks} que integrem e comparem diferentes
formas de representação e diferentes métodos de busca sobre uma mesma base de dados,
permitindo avaliar qual combinação é mais eficaz, mais interpretável e mais útil na
prática do scouting esportivo.

O presente trabalho propõe um \textit{framework} modular para busca de jogadores de
futebol a partir de estatísticas de desempenho da temporada 2024--2025, com aplicação
em tarefas de scouting, análise de mercado e comparação de atletas. O sistema compara
duas estratégias de representação textual, a concatenação de \textit{labels} e o texto
gerado por modelo de linguagem local (Llama 3.2 3B), sobre dois métodos de
recuperação: BM25 e Bi-encoder. Dois mecanismos de \textit{ground truth} garantem
a classificação confiável dos jogadores em perfis funcionais: um baseado em LLM e
outro baseado em regras ponderadas. O objetivo central é identificar qual combinação
de representação e método produz resultados ao mesmo tempo mais relevantes e mais
interpretáveis para o analista esportivo.

Além da busca por perfil, o framework incorpora um módulo de recomendação de
jogadores similares que, a partir de um jogador de referência, retorna os
atletas estatisticamente mais próximos. A similaridade é calculada por KNN com
distância euclidiana sobre um vetor de atributos normalizados por 90 minutos e
padronizados por z-score, restringindo os candidatos à mesma posição do jogador
de referência; os vizinhos são pré-computados em lote, o que torna a consulta
eficiente. Complementarmente, todas as funcionalidades são disponibilizadas por
meio de uma aplicação web interativa, que aproxima o sistema do fluxo de
trabalho do analista esportivo e permite explorar a base sem conhecimento de
programação.

Os objetivos deste trabalho organizam-se em três frentes complementares:

\begin{itemize}
  \item \textbf{Arquitetura}: construir um sistema modular para busca e recuperação
        de jogadores, com componentes independentes e configuráveis que permitam
        experimentos reprodutíveis e comparações sistemáticas entre abordagens;

  \item \textbf{Transformação}: converter os atributos estruturados de cada jogador
        em vetores numéricos e descrições textuais ricas, via concatenação de
        \textit{labels} e geração com LLM local (Llama 3.2), aplicando mecanismos
        de \textit{ground truth} (baseado em LLM e baseado em regras) para manter
        a fidelidade factual das classificações geradas;

  \item \textbf{Avaliação}: recuperar atletas por perfil técnico a partir de consultas
        em linguagem natural e validar os resultados com métricas rigorosas de IR
        (Precision@$k$, Recall@$k$ e NDCG@$k$) sobre uma estratégia
        holdout reprodutível.
\end{itemize}

O restante deste artigo está organizado da seguinte forma: a Seção~\ref{sec:fund}
apresenta a fundamentação teórica; a Seção~\ref{sec:related} discute os trabalhos
relacionados; a Seção~\ref{sec:dataset} descreve a base de dados; a
Seção~\ref{sec:dev} detalha o desenvolvimento do \textit{framework}; a
Seção~\ref{sec:exp} apresenta os experimentos e resultados; a
Seção~\ref{sec:recomendacao} descreve o módulo de recomendação de jogadores
similares; a Seção~\ref{sec:webapp} descreve a aplicação web que disponibiliza o
sistema ao usuário final; e a Seção~\ref{sec:conc} conclui o artigo.


% ============================================================
\section{Fundamentação Teórica}
\label{sec:fund}

Esta seção apresenta os conceitos e algoritmos empregados no \textit{framework},
organizados em três blocos: estratégias de representação textual dos jogadores,
mecanismos de \textit{ground truth} e controle, e estratégias de busca esparsa
(BM25) e densa (Bi-encoder), com as respectivas métricas de avaliação de ranqueamento.

\subsection{Estratégias de Representação de Jogadores}

Uma etapa central em sistemas de busca é a escolha de como representar cada item do
corpus, neste caso, cada jogador. A representação determina o espaço semântico em
que as comparações serão realizadas e influencia diretamente a qualidade dos resultados.
Neste trabalho são exploradas duas estratégias de representação textual.

\subsubsection{Label-to-Text por Concatenação}
\label{subsec:concat}

A estratégia de \textit{Label-to-Text} por concatenação transforma os atributos
estruturados de cada jogador em uma string textual organizada por categorias funcionais.
Cada categoria agrupa os atributos semanticamente relacionados (informações do jogador,
produção ofensiva, progressão, uso e disciplina), com rótulos em inglês mapeados a
partir de um dicionário de \textit{aliases} (por exemplo, \texttt{Gls}
$\rightarrow$ \texttt{goals}, \texttt{PrgC} $\rightarrow$ \texttt{progressive\_carries})
e códigos de posição expandidos para termos completos (\texttt{FW} $\rightarrow$
\texttt{forward}, \texttt{DF} $\rightarrow$ \texttt{defender}, \texttt{MF}
$\rightarrow$ \texttt{midfielder}, \texttt{GK} $\rightarrow$ \texttt{goalkeeper}).
Um exemplo de saída é:

\begin{quotation}
\textit{``player: Lionel Messi is a ARG football player who plays for Inter Miami as
a forward, midfielder (FW,MF). Age: 38. Born: 1987. Attacking output: goals 18,
assists 7, expected\_goals 12.3, expected\_assists 6.4,
non\_penalty\_expected\_goals 12.3, expected\_goal\_involvements 1.07.
Progression: progressive\_carries 90, progressive\_passes 169,
progressive\_receptions 139. Usage: matches\_played 18, starts 17,
minutes\_played 1566, full\_match\_equivalents 17.4.
Discipline: yellow\_cards 2, red\_cards 0.''}
\end{quotation}

Essa abordagem é 100\% determinística, não depende de nenhum modelo externo e produz
texto padronizado, especialmente adequado para o BM25, que se beneficia de correspondência
exata de termos. Sua principal limitação é gerar textos mecânicos que não capturam
nuances estilísticas além dos números brutos.

\subsubsection{Label-to-Text com Modelo de Linguagem Local}
\label{subsec:llm}

Esta estratégia utiliza o modelo Llama 3.2 3B executado localmente via
\texttt{llama-cpp-python} para gerar, a partir dos atributos numéricos de cada jogador,
uma descrição em português no estilo de um mini-relatório de \textit{scouting}. O
modelo é instruído por um \textit{system prompt} que o orienta a agir como analista
de futebol e produzir 3 a 4 frases factuais e objetivas, mencionando posição, clube,
nacionalidade, perfil ofensivo/defensivo e os números mais relevantes. Os parâmetros
de geração são \texttt{temperature=0.3} (baixa variabilidade) e
\texttt{max\_tokens=180}. Um exemplo de texto gerado é:

\begin{quotation}
\textit{``Lionel Messi é um jogador de futebol argentino que atua como meia-atacante
(FW/MF) e é conhecido por sua habilidade excepcional em criar oportunidades de gol.
Com 38 anos, Messi continua a ser um jogador fundamental para o Inter Miami, tendo
jogado 1566 minutos em 18 partidas, com 18 gols e 7 assistências. Ele é um especialista
em criar chances de gol, com um xG (expectativa de gol) de 12,3, e é capaz de marcar
gols e criar oportunidades para seus companheiros de equipe.''}
\end{quotation}

O \textit{system prompt} empregado instrui o modelo nos seguintes termos:

\begin{quotation}
\textit{``Você é um analista de futebol. A partir de estatísticas de um jogador,
escreva uma descrição curta (3 a 4 frases) em português, factual e objetiva.
Mencione posição, clube, nacionalidade, perfil ofensivo/defensivo e os números
mais relevantes. Não invente dados que não estão presentes.''}
\end{quotation}

Os textos gerados são armazenados em arquivo JSONL, com campos \texttt{rk},
\texttt{player}, \texttt{squad}, \texttt{pos}, \texttt{nation} e \texttt{text},
desacoplando a etapa de geração dos experimentos de recuperação. Textos gerados por
LLMs são semanticamente mais ricos e expressivos para modelos de \textit{embedding}.
Como desvantagens, o custo computacional é elevado e há variabilidade entre execuções,
já que o mesmo jogador pode receber descrições ligeiramente diferentes em chamadas
distintas ao modelo.

\subsection{Ground Truth e Controle}
\label{subsec:grounding}

Um desafio central no \textit{framework} é definir o que constitui um resultado
\textit{relevante} para cada consulta. Para isso, cada jogador recebe uma
\textbf{tag de perfil funcional}, sendo classificado em um de 12 papéis, com 3 estilos
por posição (GK, DF, MF, FW), que servem como critério de relevância nas métricas de
ranqueamento. O \textit{framework} implementa dois mecanismos alternativos para a
atribuição dessas tags.

\subsubsection{Ground Truth baseado em LLM}
\label{subsubsec:gt_llm}

Nesta abordagem, o modelo Llama 3.2 é utilizado para classificar cada jogador em uma
das 12 tags funcionais, com base exclusivamente nas suas estatísticas. O mecanismo
atua em três frentes:

\begin{itemize}
  \item \textbf{System Prompt com instruções rígidas}: o modelo é instruído a
        classificar o jogador em \textit{exatamente} uma das funções fornecidas,
        respondendo apenas com o nome da tag em \textit{snake\_case}, sem
        explicação adicional. O \textit{prompt} utilizado é:
        \begin{quotation}
          \textit{``Você é um analista de futebol. Sua tarefa é classificar um
          jogador em UMA das funções fornecidas, com base nas estatísticas. Responda
          APENAS com o nome exato da função (snake\_case), sem explicação
          adicional.''}
        \end{quotation}

  \item \textbf{Contexto restrito aos números brutos}: o \textit{user prompt} contém
        apenas as estatísticas reais do jogador, garantindo que o modelo não possa
        introduzir dados externos;

  \item \textbf{Geração determinística}: o parâmetro \texttt{temperature=0.3}
        mantém a classificação em regime de baixa criatividade, reduzindo
        inconsistências entre execuções.
\end{itemize}

A Tabela~\ref{tab:gt_llm} apresenta a distribuição das tags atribuídas pelo mecanismo
baseado em LLM sobre a base completa.

\begin{table}[!ht]
\caption{Distribuição de tags pelo Ground Truth baseado em LLM.}
\centering
\begin{tblr}{
  colspec={lcc},
  row{1}={font=\bfseries}
}
\hline
Tag & Quantidade & Porcentagem \\
\hline
zagueiro\_construtor    & 2.114 & 34,29\% \\
centroavante\_goleador  & 1.581 & 25,64\% \\
box\_to\_box            & 1.070 & 17,36\% \\
meia\_criativo          &   885 & 14,36\% \\
goleiro\_titular        &   370 &  6,00\% \\
segundo\_atacante       &    60 &  0,97\% \\
goleiro\_reserva        &    49 &  0,79\% \\
zagueiro\_classico      &    21 &  0,34\% \\
goleiro\_joia           &    14 &  0,23\% \\
volante\_defensivo      &     1 &  0,02\% \\
\hline
\end{tblr}
\label{tab:gt_llm}
\end{table}

\subsubsection{Ground Truth baseado em Regras}
\label{subsubsec:gt_rules}

Como alternativa determinística e sem dependência de modelo generativo, o
\textit{framework} implementa um mecanismo de atribuição de tags por regras
ponderadas. O processo consiste em:

\begin{itemize}
  \item \textbf{Definição dos papéis}: os mesmos 12 perfis funcionais (3 estilos
        por posição) são mantidos, garantindo comparabilidade entre os dois
        mecanismos;

  \item \textbf{Regras ponderadas}: cada papel recebe pesos sobre as colunas
        estatísticas mais relevantes para aquele perfil (por exemplo,
        \texttt{Gls} e \texttt{xG} com peso elevado para
        \texttt{centroavante\_goleador}; \texttt{PrgP} e \texttt{xAG} para
        \texttt{meia\_criativo});

  \item \textbf{Atribuição via z-score}: as estatísticas são padronizadas dentro
        de cada posição, e o \textit{score} de cada papel é calculado como
        $\text{score} = \sum_j w_j \cdot z_j$, onde $w_j$ é o peso e $z_j$ é o
        z-score do atributo $j$; o jogador recebe a tag do papel com maior
        pontuação.
\end{itemize}

A Tabela~\ref{tab:gt_rules} apresenta a distribuição das tags atribuídas pelo
mecanismo baseado em regras, que resulta em uma distribuição mais equilibrada entre
os perfis em comparação ao mecanismo baseado em LLM.

\begin{table}[!ht]
\caption{Distribuição de tags pelo Ground Truth baseado em Regras.}
\centering
\begin{tblr}{
  colspec={lcc},
  row{1}={font=\bfseries}
}
\hline
Tag & Quantidade & Porcentagem \\
\hline
zagueiro\_classico      & 1.291 & 20,94\% \\
volante\_defensivo      & 1.111 & 18,02\% \\
centroavante\_goleador  &   829 & 13,45\% \\
ponta\_driblador        &   482 &  7,82\% \\
lateral\_ofensivo       &   434 &  7,04\% \\
box\_to\_box            &   424 &  6,88\% \\
meia\_criativo          &   421 &  6,83\% \\
zagueiro\_construtor    &   410 &  6,65\% \\
segundo\_atacante       &   330 &  5,35\% \\
goleiro\_reserva        &   162 &  2,63\% \\
goleiro\_joia           &   144 &  2,34\% \\
goleiro\_titular        &   127 &  2,06\% \\
\hline
\end{tblr}
\label{tab:gt_rules}
\end{table}

\subsection{Pré-processamento e Normalização de Dados}

Antes de qualquer experimento, os dados são preparados pelas seguintes etapas:

\begin{itemize}
  \item \textbf{Seleção de atributos}: remoção de \texttt{Rk} (ranking sequencial
        sem valor preditivo);

  \item \textbf{Limpeza numérica}: conversão com \texttt{pd.to\_numeric} após remoção
        de vírgulas e ruídos textuais, com preenchimento de valores faltantes pela
        mediana de cada coluna;

  \item \textbf{Limpeza textual}: preenchimento de valores nulos com
        \texttt{``Unknown''} e remoção de espaços extras;

  \item \textbf{Deduplicação}: remoção de registros duplicados pela combinação
        \texttt{[Player, Squad]};

  \item \textbf{Normalização}: padronização \textit{z-score}
        $x' = (x - \mu) / \sigma$
        aplicada aos atributos numéricos, utilizada internamente tanto para a
        geração das estatísticas por 90 minutos nas representações textuais quanto
        para o mecanismo de Ground Truth baseado em regras.
\end{itemize}

\subsection{Busca Esparsa: BM25}

O BM25 (\textit{Best Match 25}) é um modelo probabilístico de recuperação de
informação amplamente utilizado como \textit{baseline} em sistemas de busca
\citep{robertson2009}. No contexto deste \textit{framework}, o BM25 atua sobre os
textos gerados pelas duas estratégias de representação, oferecendo alta
interpretabilidade, pois é possível inspecionar quais termos contribuíram para a
pontuação, e correspondência exata de atributos, o que o torna especialmente adequado
para o corpus de concatenação de \textit{labels}.

Dado um documento $d$ e uma consulta $q$ com termos $q_1, \ldots, q_m$, o
\textit{score} de relevância é:
\[
  \text{score}(q, d) = \sum_{i=1}^{m} \text{IDF}(q_i) \cdot
    \frac{f(q_i, d) \cdot (k_1 + 1)}
         {f(q_i, d) + k_1 \cdot \left(1 - b + b \cdot \frac{|d|}{\overline{dl}}\right)}
\]
onde $f(q_i, d)$ é a frequência do termo $q_i$ no documento $d$, $|d|$ é o comprimento
do documento, $\overline{dl}$ é o comprimento médio dos documentos, e $k_1$ e $b$ são
parâmetros de calibração. O IDF é calculado com suavização:
\[
  \text{IDF}(q_i) = \log\!\left(1 + \frac{N - df_i + 0.5}{df_i + 0.5}\right)
\]
onde $N$ é o número total de documentos e $df_i$ é o número de documentos que contêm
$q_i$.

No presente \textit{framework}, o BM25 opera sobre os corpora textuais gerados pelas
duas estratégias de representação. Os tokens são extraídos com a expressão regular
\texttt{[a-z0-9]+(?:[.+-][a-z0-9]+)*}, que preserva valores decimais como
\texttt{8.2} e termos compostos como \texttt{expected\_goals}. Os parâmetros utilizados
são $k_1 = 1.5$ (controla a saturação de frequência de termos) e $b = 0.75$
(normalização pelo comprimento do documento).

O BM25 apresenta como principais vantagens a ausência de custo de treinamento, a
interpretabilidade dos resultados e a robustez em domínios com vocabulário controlado,
como o corpus de concatenação de \textit{labels}. Como limitação, não captura
similaridade semântica entre termos distintos com mesmo significado.

\subsection{Busca Densa: Bi-encoder}

Bi-encoders são arquiteturas neurais baseadas em Transformers que codificam
separadamente uma consulta e um documento em vetores densos (\textit{embeddings})
de dimensão fixa \citep{reimers2019}. Diferentemente do BM25, que opera sobre
correspondência léxica, o Bi-encoder captura similaridade semântica profunda,
permitindo recuperar jogadores a partir de consultas em linguagem natural mesmo quando
os termos exatos não estão presentes nas descrições. A pontuação de relevância é
dada pelo produto interno entre os embeddings:
\[
  \text{score}(q, d) = E_q(q) \cdot E_d(d)
\]
Os embeddings dos documentos são pré-computados e indexados para recuperação eficiente
via \textit{Approximate Nearest Neighbor} (ANN) \citep{karpukhin2020}. No
\textit{framework} proposto, o Bi-encoder opera sobre os corpora textuais gerados
pelas estratégias de concatenação e de LLM, permitindo que consultas em linguagem
natural sejam comparadas semanticamente às descrições dos jogadores.

\paragraph{Vantagens:} captura semântica profunda; permite consultas em linguagem
natural; baixa latência na recuperação com embeddings pré-computados.

\paragraph{Desvantagens:} custo de geração dos embeddings; requer modelo pré-treinado
de qualidade; pode falhar em termos raros ou muito específicos do domínio esportivo.

\subsection{Métricas de Avaliação de Ranqueamento}

Para comparar as combinações de representação e método de recuperação, são empregadas
métricas padronizadas de avaliação de sistemas de ranqueamento \citep{manning2008}:

\begin{itemize}
  \item \textbf{Precision@$k$}: fração dos $k$ primeiros resultados que são relevantes:
        \[
          \text{P@}k = \frac{|\{\text{relevantes}\} \cap \{\text{top-}k\}|}{k}
        \]

  \item \textbf{Recall@$k$}: cobertura dos itens relevantes nos $k$ primeiros
        resultados:
        \[
          \text{Recall@}k =
          \frac{|\{\text{relevantes}\} \cap \{\text{top-}k\}|}{|\{\text{relevantes}\}|}
        \]

  \item \textbf{NDCG@$k$} (\textit{Normalized Discounted Cumulative Gain}): avalia
        a qualidade do ranqueamento considerando posição e grau de relevância:
        \[
          \text{NDCG@}k = \frac{\text{DCG@}k}{\text{IDCG@}k}, \quad
          \text{DCG@}k = \sum_{i=1}^{k} \frac{2^{rel_i} - 1}{\log_2(i+1)}
        \]
        onde $rel_i$ é o grau de relevância do resultado na posição $i$ e IDCG é o
        ganho do ranqueamento ideal.
\end{itemize}

% ============================================================
\section{Trabalhos Relacionados}
\label{sec:related}

A aplicação de aprendizado de máquina e recuperação de informação à análise e
recomendação de jogadores de futebol tem recebido atenção crescente na literatura.
Esta seção apresenta os trabalhos mais relevantes, organizados por abordagem.

\subsection{Sistemas Baseados em Similaridade Numérica}

\citet{ijai2025} propuseram um sistema de recomendação de jogadores para recrutamento
utilizando KNN com similaridade do cosseno sobre dados do FBref, com normalização por
90 minutos e redução de dimensionalidade via PCA. O modelo identificou Riyad Mahrez
como o jogador mais similar a Lionel Messi, com pontuação de 91,85\%.

\citet{researchgate2025} propuseram um \textit{framework} de scouting utilizando o
dataset do EA Sports FC 24, com busca por similaridade do cosseno e predição de
transferências via Random Forest, complementado por gráficos de radar. O sistema opera
exclusivamente sobre representações numéricas, sem explorar alternativas textuais ou
métodos de IR modernos.

\citet{playerank2019} propuseram o PlayeRank, sistema de ranqueamento validado em
dataset do Wyscout com 14 milhões de eventos, superando o estado da arte em até 16\%
na concordância com especialistas. Sua metodologia de seleção de atributos por posição
é uma referência relevante para o pré-processamento adotado neste trabalho.

\subsection{Clusterização de Jogadores por Perfil}

\citet{oconnor2023} aplicou K-Means para agrupar jogadores das cinco principais ligas
europeias por estilo de jogo, demonstrando que a clusterização pode identificar
substitutos de menor custo no mercado de transferências. \citet{carey2020} utilizou
clusterização hierárquica sobre meias do FBref, construindo uma matriz de proximidade
entre 417 jogadores sem necessidade de definir o número de grupos previamente.
\citet{aimsciences2018} aplicaram PART ao mercado do AS Roma, evidenciando que
clusterização é eficaz mesmo em dados de alta dimensionalidade.

\subsection{Recuperação de Informação Clássica e com Transformers}

O BM25 é um dos modelos de recuperação mais consolidados na literatura de IR
\citep{robertson2009} e serve como \textit{baseline} forte em benchmarks como o BEIR
\citep{thakur2021}. Sua combinação com representações textuais derivadas de dados
estruturados, como as utilizadas neste trabalho, tem sido explorada em sistemas de
busca híbridos, mas raramente aplicada ao domínio esportivo.

\citet{karpukhin2020} propuseram o DPR (\textit{Dense Passage Retrieval}), Bi-encoder
baseado em BERT treinado para recuperação de passagens em perguntas e respostas. Embora
não aplicado ao domínio esportivo, o DPR estabelece a base metodológica para o uso de
Bi-encoders em tarefas de recuperação que serão adaptadas neste trabalho.

\citet{springer2026} propuseram um sistema de recomendação combinando K-Means com Redes
Neurais Convolucionais em Grafos (GCN). Embora mais sofisticado, o método exige maior
esforço computacional e apresenta menor interpretabilidade do que as abordagens adotadas
neste trabalho.

\subsection{Posicionamento deste Trabalho}

A análise da literatura revela quatro características que diferenciam o
\textit{framework} proposto, diretamente ligadas às motivações do projeto:

\begin{enumerate}
  \item \textbf{Automatização da busca em escala}: este \textit{framework} opera sobre
        6.165 jogadores de 10 ligas simultâneas, viabilizando a recuperação automática
        de perfis em escala global a partir de consultas em linguagem natural, sem
        depender de busca manual.

  \item \textbf{Representações que tornam os dados tabulares acessíveis}: ao converter
        estatísticas brutas em texto via concatenação e LLM, o sistema transforma dados
        difíceis de interpretar em descrições consultáveis e legíveis para o analista.

  \item \textbf{Duplo mecanismo de Ground Truth}: a combinação de um mecanismo baseado
        em LLM e outro baseado em regras ponderadas para classificação de perfis
        funcionais, inédita na literatura de scouting com geração textual, permite
        avaliar a confiabilidade e o impacto do critério de relevância sobre as métricas
        de recuperação.

  \item \textbf{Avaliação por métricas padronizadas de IR com holdout}: ao contrário
        de trabalhos que usam métricas ad hoc de similaridade, este trabalho adota
        Precision@$k$, Recall@$k$ e NDCG@$k$ em uma estratégia holdout reprodutível,
        tornando os resultados comparáveis com a literatura de IR.
\end{enumerate}


% ============================================================
\section{Base de Dados}
\label{sec:dataset}

O conjunto de dados utilizado é o \textit{Top 10 Leagues Player Data 2024--25}
\citep{dataset2025}, disponível publicamente na plataforma Kaggle. A base reúne
estatísticas individuais de jogadores de futebol provenientes de dez das principais
ligas do mundo na temporada 2024--2025.

\subsection{Estrutura e Organização}

O arquivo está organizado em um documento Excel com onze abas: a aba consolidada
\textit{AllCompDataset} e dez abas específicas, uma por liga. A
Tabela~\ref{tab:ligas} apresenta a distribuição de jogadores por competição.

\begin{table}[!ht]
\caption{Distribuição de instâncias por liga na temporada 2024--2025.}
\centering
\begin{tblr}{
  colspec={lc},
  row{1}={font=\bfseries}
}
\hline
Liga & Jogadores \\
\hline
Premier League (Inglaterra)     & 574 \\
Serie A (Itália)                & 634 \\
Bundesliga (Alemanha)           & 492 \\
La Liga (Espanha)               & 601 \\
Ligue 1 (França)                & 553 \\
Brasileirão                     & 581 \\
Primeira Liga (Portugal)        & 585 \\
Belgian Pro League (Bélgica)    & 486 \\
MLS (Estados Unidos)            & 774 \\
Campeonato Argentino            & 885 \\
\hline
\textbf{Total (AllCompDataset)} & \textbf{6.165} \\
\end{tblr}
\label{tab:ligas}
\end{table}

A aba \textit{AllCompDataset} contém 6.165 instâncias e 36 atributos. Após o
pré-processamento (que inclui limpeza numérica com preenchimento pela mediana,
normalização de valores textuais e deduplicação por \texttt{[Player, Squad]})
a base não apresenta valores faltantes, garantindo que as duas estratégias de
representação partam dos mesmos dados.

\subsection{Atributos}

Os 36 atributos dividem-se em três categorias funcionais, conforme apresentado na
Tabela~\ref{tab:atributos}:

\begin{table}[!ht]
\caption{Categorização dos atributos da base de dados.}
\centering
\begin{tblr}{
  colspec={lX[l]X[l]},
  row{1}={font=\bfseries}
}
\hline
Categoria & Principais Atributos & Objetivo \\
\hline
Identificação & Nome, Nação, Clube, Idade, Posição
              & Filtros demográficos e de mercado \\
Performance   & Gols, Assistências, xG, xAG, PK
              & Avaliação de volume e eficiência ofensiva \\
Progressão    & PrgC, PrgP, PrgR, Transições
              & Métricas de avanço e construção de campo \\
\hline
\end{tblr}
\label{tab:atributos}
\end{table}

\begin{itemize}
  \item \textbf{Identificação}: \texttt{Player} (nome), \texttt{Nation} (nação),
        \texttt{Squad} (clube), \texttt{Age} (idade) e \texttt{Pos} (posição em
        campo: \texttt{FW}, \texttt{DF}, \texttt{MF}, \texttt{GK}). A base
        cobre \textbf{131 nacionalidades}, com predomínio de argentinos (915),
        brasileiros (675) e espanhóis (477);

  \item \textbf{Performance}: minutagem (\texttt{MP}, \texttt{Starts}, \texttt{Min},
        \texttt{90s}), produção ofensiva (\texttt{Gls}, máx. 39; \texttt{Ast},
        máx. 18; \texttt{G+A}; \texttt{PK}; \texttt{PKatt}) e métricas esperadas
        (\texttt{xG}, \texttt{xAG}, \texttt{npxG}), que permitem avaliar o
        desempenho além do acaso em finalizações e passes;

  \item \textbf{Progressão}: carregamentos progressivos (\texttt{PrgC}), passes
        progressivos (\texttt{PrgP}) e recebimentos progressivos (\texttt{PrgR}),
        que medem a contribuição do jogador no avanço da bola em direção ao gol
        adversário, métricas de transição de campo essenciais para identificar
        jogadores de perfil dinâmico.
\end{itemize}

Dez atributos com sufixo \texttt{.1} representam versões normalizadas por 90 minutos
das estatísticas correspondentes, permitindo comparações mais justas entre jogadores
com volumes de jogo distintos. O atributo \texttt{Rk} (ranking sequencial) é
descartado no pré-processamento por não carregar informação de desempenho.

\subsection{Natureza dos Atributos}

O conjunto combina \textbf{categóricas nominais} (\texttt{Player}, \texttt{Nation},
\texttt{Pos}, \texttt{Squad}), \textbf{numéricas discretas} (\texttt{Age}, \texttt{MP},
\texttt{Gls}, \texttt{Ast}, etc.) e \textbf{numéricas contínuas} (\texttt{90s},
\texttt{xG}, \texttt{xAG}, \texttt{npxG} e todos os atributos \texttt{.1}). A faixa
etária varia de 15 a 44 anos (média 25,4). Há 10 combinações de posição distintas,
com predomínio de defensores (\texttt{DF}, 1.801), meias (\texttt{MF}, 1.266),
atacantes (\texttt{FW}, 857) e goleiros (\texttt{GK}, 433).

\subsection{Adequação ao Problema}

A base é adequada ao problema por três razões. Primeiro, a abrangência de 6.165
jogadores de 10 ligas e 131 nacionalidades viabiliza scouting em escala global. Segundo,
as estatísticas por 90 minutos permitem comparações equilibradas entre atletas com
volumes de jogo distintos, sendo essenciais para a geração das descrições textuais.
Terceiro, a estrutura uniforme dos atributos, sem valores faltantes após o
pré-processamento, garante que as duas estratégias de representação partam de uma
base consistente.


% ============================================================
\section{Desenvolvimento do Trabalho}
\label{sec:dev}

\subsection{Arquitetura do Framework}

O \textit{framework} é implementado em Python e organizado em módulos independentes,
cada um responsável por uma etapa do \textit{pipeline}. O ponto de entrada é o
orquestrador \texttt{src/main.py}, que coordena os demais módulos: o carregamento
e pré-processamento dos dados ficam em \texttt{src/dataset/manager.py}; as estratégias
de representação textual (\texttt{concat\_labels} e \texttt{llm\_description}) são
implementadas em \texttt{src/search/dataloader/profile\_builder.py}; o motor BM25
em \texttt{src/search/model/bm25\_model.py}; e as métricas de ranqueamento em
\texttt{src/search/evaluation/}. A geração de textos via Llama e a atribuição do
\textit{ground truth} são realizadas pelo script auxiliar
\texttt{llm/scripts/textify\_players.py}, executado de forma desacoplada antes dos
experimentos de recuperação.

Os experimentos são configurados por arquivos JSON, que especificam a estratégia de
representação, o método de recuperação, os parâmetros do modelo e os caminhos de
saída, permitindo executar e comparar diferentes combinações de forma reprodutível.

\subsection{Pré-processamento e Estratégia de Amostragem}

O módulo \texttt{dataset/manager.py} realiza o carregamento e a preparação dos dados.
O arquivo Excel é lido a partir da aba \textit{AllCompDataset} (com suporte a download
automático do Kaggle quando necessário). As etapas de pré-processamento são:
(i) remoção do atributo \texttt{Rk}; (ii) limpeza e conversão numérica com
\texttt{pd.to\_numeric} após remoção de ruídos (vírgulas, caracteres estranhos),
com preenchimento de valores ausentes pela mediana de cada coluna; (iii) preenchimento
de colunas textuais com \texttt{``Unknown''} e remoção de espaços extras; e
(iv) deduplicação por \texttt{[Player, Squad]}, garantindo um registro único por
jogador por clube.

Para a avaliação dos métodos de recuperação, adota-se uma estratégia de amostragem
do tipo \textbf{holdout}: o conjunto de dados é dividido em um \textit{query set}
(conjunto de consultas de teste, com 20\% dos jogadores, amostrados aleatoriamente com
semente fixa para reprodutibilidade) e um corpus de busca (os 80\% restantes). Cada
jogador do \textit{query set} atua como consulta ao sistema, e os jogadores do corpus
com a mesma tag de perfil funcional (atribuída pelo mecanismo de \textit{ground truth}
selecionado) são considerados relevantes. Essa estratégia permite calcular as métricas
de ranqueamento de forma sistemática e comparável entre todos os métodos, sem exigir
anotações manuais de relevância.

\subsection{Geração das Representações Textuais}

As duas estratégias textuais são implementadas na classe \texttt{profile\_builder.py},
no módulo de busca.

A estratégia \texttt{concat\_labels}, implementada pela classe
\texttt{TextConcatBuilder}, converte cada linha do \textit{dataset} em uma string
organizada por categorias funcionais (informações do jogador, produção ofensiva,
progressão, uso e disciplina), usando o dicionário \texttt{LABEL\_ALIASES} para
mapear nomes técnicos de colunas para rótulos em inglês e expandindo os códigos de
posição para termos completos.

A estratégia \texttt{llm\_description}, implementada pela classe
\texttt{LLMDescriptionBuilder}, gera um texto narrativo agrupado pelas mesmas
categorias, de forma determinística, sem chamada a modelo externo. Esse texto
estruturado é mais expressivo que o formato por seções e serve tanto para o BM25
quanto para os modelos neurais de recuperação.

Para a geração real via modelo de linguagem, o script
\texttt{llm/scripts/textify\_players.py} itera sobre todos os jogadores da base e
envia seus atributos ao modelo Llama 3.2 3B, executado localmente via
\texttt{llama-cpp-python}. O \textit{ground truth} é garantido pelo \textit{system
prompt} que instrui o modelo a gerar apenas o que os números confirmam, combinado com
\texttt{temperature=0.3} para baixa variabilidade. Os textos são salvos em arquivo
JSONL em \texttt{llm/output/players\_text.jsonl}, desacoplando completamente a etapa
de geração dos experimentos de recuperação.

\subsection{Implementação dos Mecanismos de Ground Truth}

Os dois mecanismos de \textit{ground truth} são implementados como módulos
independentes e intercambiáveis no \textit{framework}.

O \textbf{Ground Truth baseado em LLM} envia as estatísticas de cada jogador ao
modelo Llama 3.2 com um \textit{system prompt} de classificação, solicitando que o
modelo retorne apenas o nome exato da tag funcional em \textit{snake\_case}. O
resultado é armazenado em arquivo JSONL, desacoplando a etapa de classificação dos
experimentos.

O \textbf{Ground Truth baseado em Regras} implementa, para cada um dos 12 perfis, um
vetor de pesos sobre as colunas estatísticas relevantes. As estatísticas são
padronizadas por z-score dentro de cada grupo de posição, e o score de cada papel é
calculado como a soma ponderada dos z-scores. O jogador recebe a tag correspondente
ao perfil de maior pontuação. Por ser totalmente determinístico e não depender de
chamadas a modelos, este mecanismo é computacionalmente eficiente e reprodutível.

\subsection{Motor de Busca BM25}

O motor BM25 é implementado do zero em \texttt{search/model/bm25\_model.py}. O
\textit{pipeline} consiste em: (i) tokenização dos documentos com a expressão regular
\texttt{[a-z0-9]+(?:[.+-][a-z0-9]+)*}, que preserva valores decimais como \texttt{8.2}
e termos compostos como \texttt{expected\_goals}; (ii) cálculo de TF por documento e
IDF com suavização; e (iii) indexação em memória dos vetores de pontuação.

Na etapa de consulta, a \textit{query} é tokenizada pelo mesmo processo, e o
\textit{score} BM25 é calculado para todos os documentos do corpus, retornando os
\texttt{top\_k} mais relevantes em um \texttt{DataFrame} com os atributos do jogador
e a pontuação obtida. Os parâmetros utilizados são $k_1 = 1.5$ e $b = 0.75$.

\subsection{Avaliação}

O módulo de avaliação (\texttt{search/evaluation/}) implementa as métricas
Precision@$k$, Recall@$k$ e NDCG@$k$. O avaliador itera sobre o \textit{query set}
do holdout, usando cada jogador como consulta, e agrega as métricas obtidas. Os
resultados são salvos em formato JSON e CSV em \texttt{outputs/results/}, permitindo
comparação direta entre as configurações de experimento.

\subsection{Estado de Implementação}

A Tabela~\ref{tab:status} resume o estado de implementação dos componentes do
\textit{framework}.

\begin{table}[!ht]
\caption{Estado de implementação dos componentes do \textit{framework}.}
\centering
\begin{tblr}{
  colspec={lc},
  row{1}={font=\bfseries}
}
\hline
Componente & Status \\
\hline
Pré-processamento completo do dataset       & Concluído \\
Geração das representações textuais         & Concluído \\
Ground Truth baseado em LLM                 & Concluído \\
Ground Truth baseado em Regras              & Concluído \\
Motor de Busca BM25                         & Concluído \\
Módulo de avaliação e métricas              & Concluído \\
Integração do Bi-encoder                    & Concluído \\
Experimentos finais e análise comparativa   & Concluído \\
\end{tblr}
\label{tab:status}
\end{table}


% ============================================================
\section{Experimentos e Análise de Resultados}
\label{sec:exp}

Esta seção apresenta os experimentos realizados com as quatro combinações de
representação e método de recuperação disponíveis (Concat+BM25, Concat+BiEncoder,
LLM+BM25 e LLM+BiEncoder), avaliadas com dois mecanismos de \textit{ground truth}
(baseado em LLM e baseado em regras) para $K \in \{10, 20, 50\}$.
Os resultados consolidados estão na Tabela~\ref{tab:tag-search-combined}.

\begin{table*}[h]
\centering
\footnotesize
\caption{Resultados da busca por perfil funcional para $K \in \{10, 20, 50\}$.
  Negrito indica o melhor valor por métrica dentro de cada bloco de $K$.}
\label{tab:tag-search-combined}
\begin{tblr}{
  colspec  = {cllcrrrrr},
  row{1}   = {font=\bfseries},
  row{2,10,18} = {bg=gray!8},
}
\hline
$K$ & Texto & Modelo & Tags & Recall & NDCG & MAP & MRR & Precision \\
\hline
\SetCell[r=8]{m} $K=10$
  & Concat & BiEncoder & LLM    & 0{,}0195          & 0{,}1658          & 0{,}0873          & 0{,}2222          & 0{,}1778 \\
  & Concat & BiEncoder & Regras & 0{,}0152          & 0{,}0953          & 0{,}0285          & 0{,}2221          & 0{,}1083 \\
  & Concat & BM25      & LLM    & 0{,}0133          & 0{,}1111          & 0{,}0503          & 0{,}2159          & 0{,}1111 \\
  & Concat & BM25      & Regras & 0{,}0080          & 0{,}0833          & 0{,}0284          & 0{,}2042          & 0{,}0833 \\
  & LLM    & BiEncoder & LLM    & 0{,}0315          & \textbf{0{,}3126} & \textbf{0{,}2529} & 0{,}3741          & \textbf{0{,}3222} \\
  & LLM    & BiEncoder & Regras & 0{,}0273          & 0{,}1789          & 0{,}0901          & 0{,}2917          & 0{,}1750 \\
  & LLM    & BM25      & LLM    & \textbf{0{,}0510} & 0{,}3126          & 0{,}2246          & \textbf{0{,}4944} & 0{,}2889 \\
  & LLM    & BM25      & Regras & 0{,}0330          & 0{,}1712          & 0{,}0838          & 0{,}3550          & 0{,}1667 \\
\hline
\SetCell[r=8]{m} $K=20$
  & Concat & BiEncoder & LLM    & 0{,}0238          & 0{,}1446          & 0{,}0589          & 0{,}2296          & 0{,}1389 \\
  & Concat & BiEncoder & Regras & 0{,}0258          & 0{,}0974          & 0{,}0260          & 0{,}2267          & 0{,}1042 \\
  & Concat & BM25      & LLM    & 0{,}0179          & 0{,}1129          & 0{,}0395          & 0{,}2159          & 0{,}1111 \\
  & Concat & BM25      & Regras & 0{,}0141          & 0{,}0833          & 0{,}0228          & 0{,}2091          & 0{,}0833 \\
  & LLM    & BiEncoder & LLM    & 0{,}0748          & \textbf{0{,}3162} & \textbf{0{,}2289} & 0{,}3741          & \textbf{0{,}3111} \\
  & LLM    & BiEncoder & Regras & 0{,}0533          & 0{,}1719          & 0{,}0676          & 0{,}3094          & 0{,}1667 \\
  & LLM    & BM25      & LLM    & \textbf{0{,}0894} & 0{,}2897          & 0{,}1733          & \textbf{0{,}4944} & 0{,}2500 \\
  & LLM    & BM25      & Regras & 0{,}0680          & 0{,}1694          & 0{,}0694          & 0{,}3648          & 0{,}1667 \\
\hline
\SetCell[r=8]{m} $K=50$
  & Concat & BiEncoder & LLM    & 0{,}0654          & 0{,}1415          & 0{,}0461          & 0{,}2336          & 0{,}1222 \\
  & Concat & BiEncoder & Regras & 0{,}0362          & 0{,}0826          & 0{,}0215          & 0{,}2267          & 0{,}0750 \\
  & Concat & BM25      & LLM    & 0{,}0314          & 0{,}1149          & 0{,}0329          & 0{,}2190          & 0{,}1111 \\
  & Concat & BM25      & Regras & 0{,}0396          & 0{,}0860          & 0{,}0170          & 0{,}2209          & 0{,}0833 \\
  & LLM    & BiEncoder & LLM    & 0{,}1556          & \textbf{0{,}3010} & \textbf{0{,}1863} & 0{,}3779          & \textbf{0{,}2689} \\
  & LLM    & BiEncoder & Regras & 0{,}1325          & 0{,}1859          & 0{,}0624          & 0{,}3117          & 0{,}1567 \\
  & LLM    & BM25      & LLM    & \textbf{0{,}2779} & 0{,}2962          & 0{,}1492          & \textbf{0{,}4976} & 0{,}2156 \\
  & LLM    & BM25      & Regras & 0{,}1833          & 0{,}2034          & 0{,}0759          & 0{,}3675          & 0{,}1650 \\
\hline
\end{tblr}
\end{table*}

\subsection{Configuração Experimental}

Os experimentos seguem a estratégia \textit{holdout} descrita na
Seção~\ref{sec:dev}: 20\% dos jogadores (1.233 instâncias, amostrados com semente
fixa) compõem o \textit{query set}; os 80\% restantes (4.932 instâncias) formam o
corpus de busca. Cada jogador do \textit{query set} é submetido como consulta ao
sistema e os jogadores do corpus com a mesma tag de perfil funcional são
considerados relevantes. As métricas Precision@$k$, Recall@$k$, NDCG@$k$, MAP e
MRR são calculadas sobre todas as consultas e reportadas como médias.

\subsection{Efeito da Representação Textual}

O fator com maior impacto nos resultados é a escolha da representação textual.
As combinações baseadas no texto gerado pelo LLM superam consistentemente as
baseadas em concatenação de \textit{labels} em todas as métricas e valores de $K$.
Para $K=10$, a melhor configuração com Concat (Concat+BiEncoder+LLM) alcança
NDCG$=0{,}1658$ e Precision$=0{,}1778$, ao passo que a melhor configuração com
texto LLM (LLM+BiEncoder+LLM) atinge NDCG$=0{,}3126$ e Precision$=0{,}3222$,
ganhos relativos de 88\% e 81\%, respectivamente. Essa vantagem se mantém estável
ao longo de todos os valores de $K$, indicando que a riqueza semântica introduzida
pelo LLM beneficia tanto o ranqueamento esparso quanto o denso.

A inferioridade do texto por concatenação era esperada para o Bi-encoder, cuja
eficácia depende de representações com vocabulário variado e frases completas. Mais
significativo é o fato de que o BM25, método projetado para correspondência léxica
exata, também melhora com o texto gerado por LLM: o BM25+Concat atinge
Recall$=0{,}0133$ em $K=10$, enquanto o BM25+LLM alcança Recall$=0{,}0510$,
melhora de 3,8$\times$. Isso sugere que a geração de texto via LLM expande o
vocabulário dos documentos com termos latentes que coincidem com os termos
naturalmente usados nas consultas, beneficiando a correspondência exata.

\subsection{Efeito do Método de Recuperação}

Comparando os dois métodos sobre o mesmo corpus textual (LLM), observa-se uma
especialização de cada abordagem. O BM25 obtém os melhores valores de Recall e MRR
em todos os $K$: em $K=50$, LLM+BM25 atinge Recall$=0{,}2779$ e MRR$=0{,}4976$,
contra Recall$=0{,}1556$ e MRR$=0{,}3779$ do LLM+BiEncoder, vantagem de 79\% e
32\%, respectivamente, nessas métricas. Isso indica que o BM25 é mais eficaz em
recuperar o primeiro resultado relevante no topo do \textit{ranking} (MRR) e em
cobrir o máximo de resultados relevantes disponíveis (Recall).

O Bi-encoder, por sua vez, apresenta vantagem clara em NDCG, MAP e Precision para
os textos LLM: em $K=10$ e $K=20$, o LLM+BiEncoder+LLM supera o LLM+BM25+LLM em
MAP (0{,}2529 vs.\ 0{,}2246 em $K=10$; 0{,}2289 vs.\ 0{,}1733 em $K=20$) e em
Precision, sinalizando que ordena os resultados relevantes de forma mais equilibrada
ao longo das $k$ posições. O BM25, por ser sensível à sobreposição léxica literal,
tende a concentrar os acertos no topo (alto MRR) mas dispersá-los ao longo do
\textit{ranking} com precisão decrescente.

Essa complementaridade sugere que uma estratégia híbrida, com BM25 para recuperação
inicial de candidatos e Bi-encoder para re-ranqueamento, pode constituir o melhor
dos dois mundos em trabalhos futuros.

\subsection{Efeito do Mecanismo de Ground Truth}

Em todas as combinações de representação e método, o \textit{ground truth} baseado
em LLM produz resultados superiores ao baseado em regras. A diferença é mais
pronunciada em MAP e Precision: para LLM+BiEncoder em $K=10$, o \textit{ground truth}
LLM alcança MAP$=0{,}2529$ contra MAP$=0{,}0901$ do baseado em regras, diferença
de 2,8$\times$. Recall e MRR apresentam diferença menor, mas ainda consistente.

Duas hipóteses explicam esse padrão. Primeiro, o \textit{ground truth} baseado em
LLM atribui tags que refletem o perfil descritivo do jogador, alinhando-se
melhor com a representação textual (também gerada por LLM): ambos descrevem o mesmo
jogador sob a mesma ótica semântica. O \textit{ground truth} por regras, por sua vez,
classifica com base em z-scores estatísticos que podem divergir da percepção funcional
capturada no texto. Segundo, a distribuição desequilibrada do \textit{ground truth}
por LLM (Tabela~\ref{tab:gt_llm}) concentra muitos jogadores em poucos perfis,
aumentando artificialmente a probabilidade de recuperar itens com a mesma tag,
o que pode inflar as métricas e deve ser considerado na interpretação dos resultados.

\subsection{Evolução das Métricas com $K$}

O Recall cresce monotonicamente com $K$ em todas as configurações, como esperado:
a melhor combinação (LLM+BM25+LLM) passa de Recall$=0{,}0510$ em $K=10$ para
$0{,}0894$ em $K=20$ e $0{,}2779$ em $K=50$, revelando que o sistema recupera
progressivamente mais jogadores relevantes ao ampliar a lista de retorno. NDCG e
Precision decrescem moderadamente com $K$ nas melhores configurações, indicando que
os resultados mais relevantes tendem a aparecer nos primeiros lugares, comportamento
desejável em sistemas de scouting, onde o analista raramente revisará mais de 20
candidatos. O MRR permanece praticamente constante com $K$ (varia menos de 0{,}5\%),
confirmando que a posição do primeiro resultado relevante é determinada pelos primeiros
resultados e não muda ao se expandir $K$.

\subsection{Síntese dos Resultados}

A análise aponta três conclusões centrais. Primeiro, a qualidade da representação
textual domina o desempenho do sistema: textos gerados por LLM superam a concatenação
de \textit{labels} em todas as métricas, independentemente do método de recuperação.
Segundo, os dois métodos de recuperação são complementares, pois o BM25 maximiza
Recall e MRR, enquanto o Bi-encoder maximiza NDCG, MAP e Precision, sugerindo que
estratégias híbridas são promissoras. Terceiro, o critério de relevância
(\textit{ground truth}) influencia significativamente as métricas, e a congruência
semântica entre o mecanismo de classificação e a estratégia de representação favorece
o \textit{ground truth} baseado em LLM. A combinação LLM+BM25+LLM é a mais robusta
para aplicações de scouting que priorizam cobertura e resposta rápida ao primeiro
resultado relevante; LLM+BiEncoder+LLM é preferível quando o objetivo é a
qualidade global do ranqueamento.

\section{Recomendação de Jogadores Similares}
\label{sec:recomendacao}

Além da busca por perfil textual, o framework disponibiliza um módulo de
\emph{recomendação de jogadores similares} que, dado um jogador de referência,
retorna os atletas estatisticamente mais próximos. Diferentemente da busca, que
parte de uma consulta em linguagem natural, a recomendação parte de um jogador
já existente na base, sendo especialmente útil para tarefas de substituição
direta, como encontrar um reposição de perfil equivalente a um atleta lesionado
ou transferido.

\subsection{Representação por Atributos por 90 Minutos}
Para tornar a comparação justa entre jogadores com volumes de jogo distintos,
o módulo utiliza exclusivamente estatísticas normalizadas por 90 minutos. O
vetor de características de cada jogador é composto por oito atributos: cinco
métricas por-90 já presentes na base (\texttt{Gls.1}, \texttt{Ast.1},
\texttt{xG.1}, \texttt{xAG.1}, \texttt{npxG.1}) e três métricas de progressão
derivadas em tempo de execução, dividindo o volume total pelo número de 90
minutos jogados (\texttt{PrgC}, \texttt{PrgP} e \texttt{PrgR} sobre \texttt{90s}).

\subsection{Padronização e Métrica de Distância}
Apenas jogadores com volume mínimo de jogo (por padrão, \texttt{90s}~$\geq 5$)
compõem o \emph{pool} qualificado. Sobre esse pool, cada atributo $j$ é
padronizado por z-score:
\begin{equation}
z_j = \frac{x_j - \mu_j}{\sigma_j},
\end{equation}
onde $\mu_j$ e $\sigma_j$ são, respectivamente, a média e o desvio-padrão do
atributo no pool qualificado (adotando-se $\sigma_j = 1$ quando nulo, para
evitar divisão por zero). A similaridade entre dois jogadores $p$ e $q$ é medida
pela distância euclidiana no espaço padronizado:
\begin{equation}
d(p, q) = \sqrt{\sum_{j=1}^{F} \left(z_j^{(p)} - z_j^{(q)}\right)^2},
\end{equation}
sendo $F = 8$ o número de atributos. Quanto menor a distância, maior a
similaridade entre o par de jogadores.

\subsection{Filtro por Posição e Seleção dos Vizinhos}
Para preservar a coerência funcional das recomendações, os candidatos são
restritos aos jogadores que compartilham a posição primária do jogador de
referência (por exemplo, um \texttt{FW} é comparado apenas com outros
\texttt{FW}), havendo a opção de desabilitar esse filtro. Para cada jogador,
calcula-se a distância a todos os candidatos e selecionam-se os $k$ vizinhos
mais próximos (KNN), com $k \in \{5, 10, 20\}$, ordenados por distância
crescente.

\subsection{Pré-computação e Armazenamento}
O cálculo dos vizinhos é realizado de forma \emph{offline}, em lote, para toda a
base. O resultado --- contendo, para cada jogador, os identificadores, a posição
e a distância dos seus \texttt{top\_5}, \texttt{top\_10} e \texttt{top\_20}
vizinhos --- é serializado em JSON e posteriormente carregado em uma tabela
relacional (\texttt{SIMILARIDADE}). Dessa forma, a recomendação em tempo de
consulta reduz-se a uma leitura indexada por jogador, garantindo latência
praticamente constante, independentemente do tamanho da base. Esse
desacoplamento entre a etapa de cálculo e a etapa de consulta espelha a
estratégia adotada para a geração das representações textuais.

\section{Aplicação Web e Interface de Usuário}
\label{sec:webapp}

Para tornar o framework acessível a analistas sem conhecimento de programação,
foi desenvolvida uma aplicação web que expõe as funcionalidades de busca,
recomendação e relatórios em uma interface interativa.

\subsection{Arquitetura da Aplicação}
A interface é implementada em \textbf{Next.js 14} (framework React com
\emph{App Router}) e \textbf{TypeScript}, com estilização em
\textbf{Tailwind CSS} e visualizações estatísticas em \textbf{Recharts}. A
persistência dos dados estruturados de jogadores, clubes, ligas e usuários
utiliza o \textbf{Supabase} (PostgreSQL gerenciado, com políticas de
\emph{Row Level Security}), acessado diretamente pelo cliente por meio de uma
camada de serviços, organizada em um módulo por entidade (jogadores,
estatísticas, clubes, ligas, nacionalidades, pesquisas e usuários).

A busca por perfil, por sua vez, é atendida por um microsserviço em Python que
expõe os motores de recuperação via HTTP. A aplicação web encaminha as consultas
a esse serviço por meio de um \emph{proxy} de rota (\texttt{/api/ml/*}),
mantendo o motor de recuperação desacoplado da camada de apresentação. Essa
separação permite evoluir os algoritmos de IR (BM25 e Bi-encoder) sem impacto na
interface.

\subsection{Funcionalidades}
A aplicação organiza-se nas seguintes telas:
\begin{itemize}
\item \textbf{Painel (Overview)}: indicadores gerais da base (total de
jogadores, times, ligas e nacionalidades), distribuição por posição e destaques
de artilheiros e assistentes;
\item \textbf{Jogadores}: listagem com filtros combinados por nome, posição,
time e liga, com autocompletar e paginação;
\item \textbf{Perfil do Jogador}: página detalhada com estatísticas por
temporada e gráficos que contrastam desempenho realizado e esperado (xG/xAG),
além das ações progressivas;
\item \textbf{Busca por Perfil}: consulta em linguagem natural atendida pelos
motores de recuperação;
\item \textbf{Recomendar Similares}: recomendação de jogadores estatisticamente
próximos, conforme descrito na Seção~\ref{sec:recomendacao};
\item \textbf{Relatórios}: geração de rankings por métrica e de comparativos de
similaridade, com exportação em PDF;
\item \textbf{Autenticação e Histórico}: cadastro e login de usuários e registro
das últimas buscas por tela, persistidas no banco (usuário autenticado) ou
localmente no navegador.
\end{itemize}

\subsection{Seleção do Motor de Recuperação}
Na tela de busca por perfil, o usuário pode selecionar qual motor de recuperação
será utilizado para atender à consulta, alternando entre \textbf{BM25} (busca
esparsa, por correspondência léxica) e \textbf{Bi-encoder} (busca densa, por
similaridade semântica). A escolha é enviada como parâmetro ao microsserviço, que
roteia a consulta ao motor correspondente e retorna os \texttt{top\_k} jogadores
mais relevantes. Essa flexibilidade permite ao analista contrastar, na prática, o
comportamento das duas estratégias descritas nas Seções~2.4 e~2.5 e comparadas
experimentalmente na Seção~6: o BM25, mais eficaz em cobertura e em posicionar
rapidamente o primeiro resultado relevante, e o Bi-encoder, mais preciso na
qualidade global do ranqueamento.

% ============================================================
\section{Conclusão}
\label{sec:conc}

Este trabalho apresentou um \textit{framework} modular para busca de jogadores de
futebol a partir de estatísticas de desempenho da temporada 2024--2025, com foco em
aplicações de scouting, análise de mercado e comparação de atletas. O sistema integra
e compara duas estratégias de representação textual (concatenação de \textit{labels}
e geração via LLM local, Llama 3.2 3B) sobre dois métodos de recuperação (BM25 e
Bi-encoder), avaliados com dois mecanismos de \textit{ground truth} (baseado em LLM
e baseado em regras) e métricas padronizadas de IR em uma estratégia \textit{holdout}
reprodutível.

Os experimentos revelaram que a escolha da representação textual é o fator mais
determinante do desempenho: textos gerados por LLM superam a concatenação de
\textit{labels} em até 88\% em NDCG e em até 81\% em Precision em $K=10$,
beneficiando tanto o BM25 quanto o Bi-encoder. Os dois métodos de recuperação
mostraram-se complementares: o BM25 maximiza Recall e MRR, recuperando mais
jogadores relevantes e posicionando o primeiro resultado relevante mais alto no
\textit{ranking}; o Bi-encoder maximiza NDCG, MAP e Precision, ordenando os
resultados de forma mais equilibrada. O mecanismo de \textit{ground truth} baseado
em LLM produziu resultados consistentemente superiores ao baseado em regras, o que
é parcialmente explicado pela congruência semântica entre a classificação e a
representação textual, ambas geradas pela mesma base de modelo de linguagem.

A arquitetura modular do \textit{framework}, com configuração por arquivos JSON,
separação entre geração de representações e experimentos de recuperação, e
intercambialidade dos mecanismos de \textit{ground truth}, facilita a reprodução
e a extensão dos experimentos. A estratégia \textit{holdout} com semente fixa garante
que as comparações entre combinações sejam justas e comparáveis.

Para além da avaliação dos métodos de recuperação, o framework foi estendido com
um módulo de recomendação de jogadores similares e com uma aplicação web. O
módulo de recomendação identifica, para um jogador de referência, os atletas
mais próximos por meio de KNN com distância euclidiana sobre um vetor de
atributos por 90 minutos padronizados por z-score, restringindo os candidatos à
mesma posição e pré-computando os vizinhos em lote para consulta eficiente. A
aplicação web, por sua vez, disponibiliza a busca por perfil, com escolha
entre os motores BM25 e Bi-encoder, a recomendação de similares e a geração
de relatórios em uma interface interativa, aproximando o sistema do uso prático
por analistas esportivos.

Do ponto de vista prático, os resultados recomendam a combinação LLM+BM25 para
cenários em que o analista precisa cobrir rapidamente o maior número de candidatos
relevantes, e LLM+BiEncoder para cenários em que a precisão do ranqueamento é
prioritária. Uma estratégia híbrida que utilize BM25 para recuperação inicial de
candidatos e Bi-encoder para re-ranqueamento pode reunir as vantagens de ambos e
constitui a principal direção para trabalhos futuros.

Como limitações, destacam-se: (i) a dependência do custo computacional da geração
dos textos via LLM local, que inviabiliza atualizações frequentes da base sem
infraestrutura adequada; (ii) o potencial viés introduzido pelo desequilíbrio na
distribuição das tags do \textit{ground truth} baseado em LLM, que pode inflar as
métricas de recuperação; e (iii) a cobertura restrita ao corpus da temporada
2024--2025, que não captura trajetórias longitudinais dos atletas.

Como trabalhos futuros, pretende-se: (i) implementar a estratégia híbrida
BM25+Bi-encoder com re-ranqueamento; (ii) investigar modelos de \textit{embedding}
especializados em domínio esportivo; (iii) avaliar a aceitação do sistema por
analistas de futebol em estudos de usabilidade; e (iv) estender a base para temporadas
anteriores, permitindo análise de trajetória e predição de desempenho futuro.

% ============================================================
\begin{declarations}

\begin{materials}
O conjunto de dados está disponível publicamente em:
\url{https://www.kaggle.com/datasets/abhayr10/top-10-leagues-player-data2024-25}.
Os códigos fonte desenvolvidos estão disponibilizados em repositório público no GitHub,
no seguinte link:
\url{https://github.com/MessiasFCM/Framework-football_player} .
\end{materials}

\end{declarations}

% ============================================================
\bibliographystyle{apalike-sol}
\bibliography{refs}

\end{document}